\scanpage{136}
% archive sheet number: 135
\diagram{\resizebox{\linewidth}{!}{\begin{tikzcd}[ampersand replacement=\&,column sep=small,row sep=small]
\& \operatorname{Pro}(\mathrm{Sch}_{/S}) \& \text{pro-préschémas}\\
\& \bigcup\ \substack{\text{sous-catégorie}\\\text{pleine}} \&\\
(\mathrm{Sch}_{/S}) \& \operatorname{Pro}'(\mathrm{Sch}_{/S})\arrow[l,"\varprojlim"'] \& \text{pro-préschémas} \& (X_i)_{i\in I},\quad\text{t.q.}\quad\exists i,\quad u_{ji}\ \text{affine pour }j\geq i\\
\& \bigcup\ \substack{\text{sous-catégorie}\\\text{pleine}} \&\\
\& \operatorname{Pro}''(\mathrm{Sch}_{/S})\arrow[uul,"\substack{\text{foncteur}\\\text{\unsure{plein\textsuperscript{mt}|plein|pleinement}}\\\text{fidèle}}"'] \& \text{pro-préschémas} \& {\begin{aligned}
&(X_i)_{i\in I}\quad\text{t.q.}\quad
\deleted{\left\{\begin{aligned}
&\unsure{X_i\ \text{présentation finie}}\\[-2pt]
&\exists i,\ \bigl[\unsure{X_i\ \text{quasi cpt}|X_i\ \text{qu.-compact et}}\ \text{de prés. finie}\\
&\quad u_{ji}\ \text{affine de prés. finie}
\end{aligned}\right.}\\
&\left\{\begin{aligned}
&\forall i\ \unsure{\mathrm{grd}|\mathrm{grand}},\ X_i\ \text{est qu. cpt et de présentation finie sur }S\\
&\exists i\ \text{t.q.}\ j\geq i\ \text{implique }u_{ji}\ \text{affine.}
\end{aligned}\right.
\end{aligned}}
\end{tikzcd}}}

Soient $\mathscr X\in\operatorname{Ob}\operatorname{Pro}''(\mathrm{Sch}_{/S})$, $\mathscr Y\in\operatorname{Ob}\operatorname{Pro}'(\mathrm{Sch}_{/S})$ \inserted{$\mathscr Y=(Y_j)$, les $Y_j$ \unsure{qu.cpts|quasi cpts},} alors
\[
\operatorname{Hom}_{\operatorname{Pro}(\mathrm{Sch}_{/S})}(\mathscr Y,\mathscr X)
\xrightarrow{\sim}\operatorname{Hom}(\varprojlim\mathscr Y,\varprojlim\mathscr X)
\]

\deleted{[En particulier, la catégorie}\inserted{Si \unsure{on|tout} est sur une base $S$}
\[
\operatorname{Hom}_{\operatorname{Pro}(\mathrm{Sch}_{/S})}(\mathscr Y,\mathscr X)
\simeq\operatorname{Hom}_S(\varprojlim\mathscr Y,\varprojlim\mathscr X)
\]

\deleted{Dans les provariétés de Serre}

\uline{Remarque}\quad Dans Serre et Néron, on utilise des systèmes projectifs de schémas \inserted{(algébriques)} qui sont \inserted{(pratiquement tjrs)} $\in\operatorname{Pro}''(\mathrm{Sch})$. \deleted{\unsure{Du|D}} Il y a avantage à les ``interpréter'' comme étant des véritables schémas sur $k$, \deleted{\unsure{de}} au lieu de pro-objets \ldots\ldots
